\section{A uniform frequency-block estimate}\label{sec:frequency}

The rough kernels arising from hard truncation will be decomposed into
frequency blocks.  The estimate in this section permits their coefficients
to depend on the output point and on frequency.  Its decisive hypothesis is
a pointwise square-sum bound over scales and frequencies; no sum of
frequency suprema is required.

Fix $q>0$, let $D\geq1$, and let $\Xi_D$ be a measurable subset of
$[-2D,2D]$, equipped with the measure $\dd\gamma(\xi)=\dd\xi/D$.
Thus $\gamma(\Xi_D)\leq4$.  Define
\begin{equation}\label{eq:frequency-windows}
 \ell_\xi(z)=D^{-1}\frac{\dd}{\dd z}
                 \bigl(g_q(z)e^{i\xi z}\bigr),
 \qquad H_\xi=\ell_\xi*\ell_\xi*\ell_\xi.
\end{equation}
The parameters $q,D$ are suppressed in this notation.  Constants below
may depend on $q$, but not on $D$.

\begin{proposition}[Frequency-block estimate]\label{prop:frequency-block}
Let $\mathcal I\subset\mathbb Z$ be finite and nonempty, and set
$L_k=2^k$.  Let $A_0,A_1,A_2$ be finite arrays on the lattice
$u_i=ih$, with the norms $n_v$ defined in the heat-flow preliminaries.
Suppose that measurable coefficients $\mu_{ij,k}:\Xi_D\to\C$ satisfy
\begin{equation}\label{eq:frequency-coefficient-condition}
 |\mu_{ij,k}(\xi)|\leq C_1,
 \qquad
 \sum_{k\in\mathcal I}\int_{\Xi_D}
           |\mu_{ij,k}(\xi)|^2\,\dd\gamma(\xi)\leq C_1^2
\end{equation}
for every output index pair $(i,j)$, with the first inequality holding
almost everywhere.  If
\begin{equation}\label{eq:frequency-mesh}
 h^2\leq\frac{q\min_{k\in\mathcal I}L_k^2}{16D^2},
\end{equation}
then
\begin{equation}\label{eq:frequency-finite-estimate}
 \left|h^3\int_{\Xi_D}\sum_{k\in\mathcal I}\sum_{i,j,l}
 A_0(i,j)\mu_{ij,k}(\xi)A_1(j,l)A_2(l,i)
 \Dil_{L_k}H_\xi(u_i+u_j+u_l)\,\dd\gamma(\xi)\right|
 \leq C(q,C_1)\prod_{v=0}^2n_v.
\end{equation}
\end{proposition}

There are two reasons for refining one Gaussian variance in the proof.
First, a window oscillating at frequency $D$ becomes an average of score
insertions at variance comparable to $D^{-2}s$.  Second, the boundary
energy produced by changing the output array can still be estimated
uniformly at this small variance.  More precisely, the refined row
scores will be controlled by the heat identity through single-edge
energies and the changes of star energy at slab boundaries.  Such a
change contains the square of the changing edge norm times the norm
of the unchanged edge.  We isolate a uniform estimate for that mixed
term before beginning the main proof.

\begin{lemma}[A uniform mixed boundary estimate]
\label{lem:frequency-boundary}
Fix distinct vertices $v,w\in\{0,1\}$, so that $\{v,w\}=\{0,1\}$.
Let the edge joining $v,w$ have array $B$, indexed in the original
endpoint order $(0,1)$, and let the edge joining $v,2$ have array $E$,
indexed in the order $(v,2)$.  Use variances
$(\sigma_v,\sigma_w,\sigma_2)=(\theta s,s,s)$, where
$0<\theta\leq1/16$ and $h^2\leq\theta s$.  Write
$\dd\nu_s(d)=g_{(1-\theta)s}(d)\,\dd d$.
There is a nonnegative array $\mathcal H_E(i,j)$, independent of
$s,\theta$, such that
\begin{align}
 &\int_{\R}\int_{\Pi_d}
       \|C_{v,w}\|_{\HS}^2\|C_{v,2}\|_{\HS}
                  \,\dd\pi_d\,\dd\nu_s(d)
 \leq C\sum_{i,j}h^2|B(i,j)|^2\mathcal H_E(i,j),
 \label{eq:frequency-boundary-mixed}\\
 &\left(\sum_{i,j}h^2\mathcal H_E(i,j)^3\right)^{1/3}
 \leq C\left(\sum_{m,l}h^2|E(m,l)|^3\right)^{1/3}.
 \label{eq:frequency-boundary-norm}
\end{align}
Here $(i,j)$ are always in the original order of the endpoints of
edge $0$; $i_v,i_w$ denote the same two indices in the order $(v,w)$.
The input arrays $B,E$ are extended by zero to $\mathbb Z^2$.
\end{lemma}

\begin{proof}
Expand the square of the first Hilbert--Schmidt norm.  For each pair
$(i,j)$, the measure
\[
 h^{-2}w_v(i_v)w_w(i_w)\,\dd\pi_d\,\dd\nu_s(d)
\]
is a probability measure.  Indeed, with $p_v,p_w$ as the free
coordinates on $\Pi_d$, it makes $p_v,p_w,d$ independent normal
variables with respective means $u_{i_v},u_{i_w},0$ and variances
$\theta s,s,(1-\theta)s$.  The remaining center is
$p_2=d-p_v-p_w$.

Under this probability measure, the expectation of
$h^{-2}w_v(m)w_2(l)$ is a bivariate Gaussian density, evaluated at
$(u_m,u_l)$.  Its mean is
$(u_{i_v},-u_{i_v}-u_{i_w})$, and its covariance matrix is
\[
 s\begin{pmatrix}2\theta&-\theta\\-\theta&3\end{pmatrix}.
\]
Call this covariance matrix $\Sigma$.  The bounds
$\Sigma\leq4s\operatorname{diag}(\theta,1)$ and
$\det\Sigma=\theta(6-\theta)s^2$ show, directly from the Gaussian
density formula, that
\begin{equation}\label{eq:frequency-covariance-domination}
 \mathbb E[w_v(m)w_2(l)]
 \leq Ch^2g_{C'\theta s}(u_m-u_{i_v})
                g_{C's}(u_l+u_{i_v}+u_{i_w}),
\end{equation}
where $C$ is absolute and one may take $C'=4$.

Let $M_{\mathbb Z}^{(1)}$ and $M_{\mathbb Z}^{(2)}$ be the discrete
line maximal operators in the two coordinates.  By
Lemma~\ref{lem:line-maximal}, we can take
\begin{equation}\label{eq:frequency-boundary-majorant}
 \mathcal H_E(i,j)^2
 =C\bigl(M_{\mathbb Z}^{(1)}M_{\mathbb Z}^{(2)}|E|^2\bigr)
                 (i_v,-i_v-i_w).
\end{equation}
In fact, \eqref{eq:frequency-covariance-domination} and
$h^2\leq\theta s$ show that
$\mathbb E\|C_{v,2}\|_{\HS}^2\leq\mathcal H_E(i,j)^2$.
Cauchy--Schwarz under this probability measure now proves
\eqref{eq:frequency-boundary-mixed}.  The map
$(i,j)\mapsto(i_v,-i_v-i_w)$ is a bijection of $\mathbb Z^2$.
The $\ell^{3/2}$ bounds for the two line maximal operators, applied
to $|E|^2$, prove \eqref{eq:frequency-boundary-norm}.
\end{proof}

\begin{proof}[Proof of Proposition~\ref{prop:frequency-block}]
We may discard a common null set of frequencies so that all pointwise
bounds in \eqref{eq:frequency-coefficient-condition} hold.  Set
\begin{equation}\label{eq:frequency-slabs}
 B_{0,k}^\xi(i,j)=A_0(i,j)\mu_{ij,k}(\xi),\qquad
 a_k=qL_k^2,\qquad
 \theta=\frac1{16D^2}.
\end{equation}
For each $\xi$, use $B_{0,k}^\xi$ on the open slab
$2a_k<s<3a_k$, and use the zero array on edge $0$ outside these
slabs.  The slabs are pairwise disjoint because $a_{k+1}=4a_k$.
The other two arrays remain $A_1,A_2$ throughout.  All scale
integrals below are over
\[
 [s_-,s_+]=[\min_{k\in\mathcal I}a_k,
                         4\max_{k\in\mathcal I}a_k].
\]
The mesh hypothesis gives $h^2\leq\theta s$ on this entire interval.

We first turn the oscillatory windows into refined row scores.  We
then bound the integrated row scores by telescoping the star energies,
using Lemma~\ref{lem:frequency-boundary} at every slab boundary.

\paragraph{Refining a row variance.}
Fix a slab, put $L=L_k$, $a=a_k$, and $\omega=\xi/L$.  The window
at each vertex is
\begin{equation}\label{eq:frequency-scaled-window}
 f_{\xi,k}(u):=\Dil_L\ell_\xi(u)
   =\frac LD\frac{\dd}{\dd u}
                  \bigl(g_a(u)e^{i\omega u}\bigr).
\end{equation}
Initially give all three vertices the variance $s$.  The diagonal
insertion with entries $f_{\xi,k}(u_i-p_v)/g_s(u_i-p_v)$ has
operator norm at most $C(q)$, uniformly on $2a<s<3a$.  Indeed,
$|L\omega/D|\leq2$, and the remaining factor is bounded using
Gaussian decay in $g_a/g_s$.  Split this diagonal into its real
and imaginary parts before applying Lemma~\ref{lem:mixed-trace}.
The resulting eight terms involve only Hermitian insertions of
uniformly bounded norm.

Consider the cost belonging to the row star $R_v$.  Replace its row
variance $s$ by $b'=\theta s$, leave both column variances equal
to $s$, and shift the row center from $p_v$ to $p_v+d$.  Denote
the refined Gram and score insertion by $T_v^d$ and $U_v^d$.
Addition of Gaussian variances gives
\begin{equation}\label{eq:frequency-gram-average}
 T_v=\int_{\R}T_v^d\,\dd\nu_s(d),
 \qquad \dd\nu_s(d)=g_{s-b'}(d)\,\dd d.
\end{equation}
All these Grams have a common support: changing a strictly positive
row density amounts to invertible row scaling and leaves the kernel
of the row matrix unchanged.

To represent the inserted Gram in the same way, put
$\omega'=a\omega/(a-b')$.  Since $b'<a/2$, completing the square
gives the exact identity
\begin{equation}\label{eq:frequency-modulated-gaussian}
 g_a(u)e^{i\omega u}
 =\exp\!\left(\frac{\omega^2ab'}{2(a-b')}\right)
       \int_{\R}g_{b'}(u-d)e^{i\omega'd}g_{a-b'}(d)\,\dd d.
\end{equation}
The exponential prefactor is at most $C(q)$, because
$\omega^2b'\leq C(q)$.  Also
\[
 g_{a-b'}(d)\leq Cg_{s-b'}(d),\qquad
 \frac LD\leq C(q)\sqrt{\theta s}.
\]
Differentiate \eqref{eq:frequency-modulated-gaussian} in $u$.
The derivative of the first Gaussian supplies the negative refined
row score $-(u-d)/b'$.  Consequently, for either the real or the
imaginary diagonal insertion, its inserted Gram $K_v$ has the form
\[
 K_v=\int_{\R}\alpha(d)U_v^d\,\dd\nu_s(d),
 \qquad \alpha(d)=\operatorname{Re}c(d)
       \quad\hbox{or}\quad\operatorname{Im}c(d),
\]
where the scalar coefficient is explicitly
\[
 c(d)=-\frac LD
       \exp\!\left(\frac{\omega^2ab'}{2(a-b')}\right)
       e^{i\omega'd}\frac{g_{a-b'}(d)}{g_{s-b'}(d)}.
\]
The preceding bounds give
$|\alpha(d)|\leq C(q)\sqrt{\theta s}$.
The averaged matrices are integrable, since the finite-array Gaussian
weights dominate the polynomial score factors.  By
\eqref{eq:frequency-gram-average},
Lemma~\ref{lem:joint-convexity}, and the quadratic homogeneity of
$b_T$ in its second argument,
\begin{equation}\label{eq:frequency-score-refinement}
 b_{T_v}(K_v)
 \leq C(q)\theta s\int_{\R}
           V_v(p_0,\ldots,p_v+d,\ldots,p_2;s)\,\dd\nu_s(d).
\end{equation}
The cost on the right uses variances
\begin{equation}\label{eq:frequency-anisotropic-variances}
 \sigma_v=\theta s,\qquad \sigma_w=s\quad(w\ne v).
\end{equation}
Thus the factor $\theta$ in the refined score is gained at exactly
the scale needed to accommodate the oscillation.

\paragraph{Integrated refined scores.}
Fix $v$ and use \eqref{eq:frequency-anisotropic-variances} for the
rest of this part of the proof.  All stars, energies, and costs
refer to the slab-dependent edge $0$ just defined.  Dependence
on $\xi$ is suppressed until frequency integration is needed.
For a center-dependent cost $Q$, write
\[
 \widetilde Q(s)=\int_{\R}\int_{\Pi_d}Q(p,s)
                         \,\dd\pi_d\,\dd\nu_s(d),\qquad
 \widetilde J_v(s)=\int_{\R}J_v(d,s)\,\dd\nu_s(d).
\]
On each open interval where the arrays are fixed, the heat identities
\eqref{eq:heat-plane-identities} give
\begin{equation}\label{eq:frequency-averaged-heat}
 2\partial_s\widetilde J_v
   =3\widetilde Z_v-\theta\widetilde V_v
                           -\sum_{w\ne v}\widetilde Y_{v,w}.
\end{equation}
Here the coefficient $3$ has a useful interpretation.  The original
plane second derivative has coefficient $2+\theta$.  Differentiating
$\nu_s$, applying its heat equation, and integrating twice by parts
adds $1-\theta$.  The second heat identity identifies the averaged
plane second derivative with $\widetilde Z_v$.

For fixed arrays on a compact positive scale interval,
Lemma~\ref{lem:heat-identities} bounds the plane integrals of the
energies, costs, and relevant absolute derivatives by polynomials in
$|d|$. The density $\nu_s$ and its derivatives have Gaussian decay.
These bounds justify differentiation under the integral and the two
integrations by parts in $d$, and give the required one-sided limits
at slab boundaries. Their constants may depend on the mesh and
$\theta$; they justify the identities and do not enter the estimates.

The metric Cauchy--Schwarz inequality gives
$2Z_v\leq\sum_{w\ne v}Y_{v,w}$.  Hence
\begin{equation}\label{eq:frequency-score-differential}
 \theta\widetilde V_v
 \leq-2\partial_s\widetilde J_v
                         +\frac12\sum_{w\ne v}\widetilde Y_{v,w}.
\end{equation}
The column costs in this inequality can be paid for by individual
edges.  Indeed, $S_v\geq C_{v,w}C_{v,w}^*$, so the order comparison
of Lemma~\ref{lem:order} yields
\begin{equation}\label{eq:frequency-edge-cost}
 Y_{v,w}\leq
 b_{C_{v,w}C_{v,w}^*}(C_{v,w}N_wC_{v,w}^*)=:Y^e_{v,w}.
\end{equation}
The inserted matrix is supported on the single-edge Gram, so this
comparison also applies if the larger star Gram has a larger
support.  Define the integrated single-edge energy
\[
 J^e_{v,w}(s)=\iint_{\R^2}
                  \Phi(C_{v,w}C_{v,w}^*)\,\dd p_v\,\dd p_w.
\]
The integral of $Y^e_{v,w}$ over $\Pi_d$ is independent of $d$,
because only the two endpoint centers occur.  With just this edge
present, the heat identities give
\begin{equation}\label{eq:frequency-edge-dissipation}
 -2\partial_sJ^e_{v,w}(s)
 \geq\iint_{\R^2}Y^e_{v,w}(p,s)\,\dd p_v\,\dd p_w.
\end{equation}
The coefficient of this column cost is $1$; the remaining row cost
is nonnegative.  Thus the estimate is independent of the small
variance ratio $\theta$.

If the edge is unchanged throughout $[s_-,s_+]$, integration of
\eqref{eq:frequency-edge-dissipation} and
Lemma~\ref{lem:single-edge-energy} bound its total cost by
$C(n_1^3+n_2^3)$.  If it is edge $0$, apply the same argument
separately on its slabs.  The result is bounded by
\[
 C\sum_{k\in\mathcal I}\sum_{i,j}h^2
                         |B_{0,k}^\xi(i,j)|^3.
\]
After integration in $\xi$, this is at most $C(C_1)n_0^3$, since
\begin{equation}\label{eq:frequency-cubic-summation}
 \sum_k\int_{\Xi_D}|B_{0,k}^\xi(i,j)|^3\,\dd\gamma
 \leq C_1|A_0(i,j)|^3
       \sum_k\int_{\Xi_D}|\mu_{ij,k}(\xi)|^2\,\dd\gamma
 \leq C_1^3|A_0(i,j)|^3.
\end{equation}
We have therefore controlled all column terms in
\eqref{eq:frequency-score-differential}.  The remaining issue is
the boundary energy of the stars.

\paragraph{Summing the slab boundaries.}
Before the first slab, edge $0$ is zero.  By
Lemma~\ref{lem:initial-energy},
$\widetilde J_v(s_-)\leq C(n_1^3+n_2^3)$.
The star at vertex $2$ has no jumps.  For $v\in\{0,1\}$, let
$w$ be the other vertex of edge $0$.  At either boundary of a
slab, write $C_{\rm sw}=C_{v,w}$ for the slab-side weighted block and
$C_{\rm fix}=C_{v,2}$ for the unchanged weighted block.
The absolute change in the pointwise star
energy is at most
\begin{equation}\label{eq:frequency-pointwise-jump}
 C\bigl(\|C_{\rm sw}\|_{\HS}^3
       +\|C_{\rm sw}\|_{\HS}^2\|C_{\rm fix}\|_{\HS}\bigr).
\end{equation}
To verify this bound, integrate the derivative of
$\Phi(C_{\rm fix}C_{\rm fix}^*+tC_{\rm sw}C_{\rm sw}^*)$
for $0<t\leq1$. Its support is fixed there,
and the derivative is
\[
 \begin{aligned}
 &\frac32\tr\bigl((C_{\rm fix}C_{\rm fix}^*
                   +tC_{\rm sw}C_{\rm sw}^*)^{1/2}
                    C_{\rm sw}C_{\rm sw}^*\bigr)\\
 &\qquad\leq\frac32\bigl(\|C_{\rm fix}\|_{\HS}
                    +\|C_{\rm sw}\|_{\HS}\bigr)
                      \|C_{\rm sw}\|_{\HS}^2.
 \end{aligned}
\]
Continuity gives the endpoint at $t=0$.  Removing the edge has
the same absolute change as adding it.

The integrated cubic term in
\eqref{eq:frequency-pointwise-jump} is bounded by the cube of
the norm of $B_{0,k}^\xi$, by
Lemma~\ref{lem:initial-energy}.  Summing both boundaries of
every slab and using \eqref{eq:frequency-cubic-summation} gives
$C(C_1)n_0^3$ after frequency integration.

For the mixed term, let $A^{(v,2)}$ be the fixed unweighted array on
the edge joining $v,2$, indexed in that order. Apply
Lemma~\ref{lem:frequency-boundary} with this array and denote its
majorant by $\mathcal H_v$. It is independent of $k,\xi,s$, and its cube norm is
at most $C(n_1+n_2)$.  Both boundaries together consequently
contribute at most a constant times
\begin{align*}
 &\sum_{i,j}h^2|A_0(i,j)|^2\mathcal H_v(i,j)
       \sum_k\int_{\Xi_D}|\mu_{ij,k}(\xi)|^2\,\dd\gamma\\
 &\hspace{25mm}\leq C(C_1)n_0^2(n_1+n_2),
\end{align*}
by \eqref{eq:frequency-coefficient-condition} and H\"older's
inequality.  This is the point at which averaging coefficient
squares over frequency is essential.  Taking a separate supremum
in frequency for each scale would not give this bound.

Integrate \eqref{eq:frequency-score-differential} on the intervals
between consecutive boundaries and telescope.  The final energy
is nonnegative and may be discarded.  The initial energy and
the absolute jumps just estimated, together with the single-edge
costs, give
\begin{equation}\label{eq:frequency-integrated-score}
 \int_{\Xi_D}\int_{s_-}^{s_+}
        \theta\widetilde V_v(s)\,\dd s\,\dd\gamma(\xi)
 \leq C(C_1)\sum_{j=0}^2n_j^3.
\end{equation}
Here we absorbed $n_0^2(n_1+n_2)$ by Young's inequality.
Values at the finitely many boundaries are immaterial.
All integrands are measurable in $\xi$: the arrays are measurable,
and the finite-matrix energies and supported quadratic forms are
Borel functions of their entries.  The latter can also be obtained
as limits of positive-definite regularizations.  Thus every
nonnegative frequency integration above is justified by Tonelli's
Theorem, without any regularity assumption on the coefficient
functions beyond measurability.

\paragraph{Recovering the frequency-block form.}
Return to the equal-variance stars and apply
Lemma~\ref{lem:mixed-trace} to the real and imaginary parts of
the three window insertions.  Integrate the bound over base centers
on $\Pi_0$.  Shifting the row center in
\eqref{eq:frequency-score-refinement} maps this plane to $\Pi_d$;
using the two column centers as free coordinates shows that its
Jacobian is $1$.  The integrated cost at vertex $v$ is therefore
at most
$C(q)\theta s\widetilde V_v(s)$.
Integrate over every slab with $\dd s/s$, and then over $\Xi_D$
with $\dd\gamma$.  The sum of the resulting bounds is controlled
by \eqref{eq:frequency-integrated-score}.

On the other side, multiplying the three weighted windows in
the trace and integrating the centers gives exactly
\[
 h^3\sum_{i,j,l}B_{0,k}^\xi(i,j)A_1(j,l)A_2(l,i)
                \Dil_{L_k}H_\xi(u_i+u_j+u_l).
\]
This is the convolution identity on $\Pi_0$: the three arguments
of the windows sum to $u_i+u_j+u_l$.  The expression is independent
of $s$ within its slab, whose logarithmic length is
$\int_{2a_k}^{3a_k}\dd s/s=\log(3/2)$.
All signed center integrations are absolutely convergent by
Gaussian decay.  We have proved
\eqref{eq:frequency-finite-estimate} with its right side replaced
by $C(q,C_1)\sum_v n_v^3$.

Finally, if all $n_v$ are nonzero, apply this estimate to
$A_v/n_v$.  The coefficients and mesh hypothesis are unchanged,
and trilinearity gives the product bound claimed in the
Proposition.  If any norm is zero, the form is zero.
\end{proof}

\begin{corollary}[Continuous frequency-block estimate]
\label{cor:frequency-continuum}
Let $q,D,\Xi_D,\mathcal I$ be as in
Proposition~\ref{prop:frequency-block}.  Suppose that
$\mu_{a,b,k}(\xi)$ is constant in $(a,b)$ on every cell of one
finite rectangular partition, independent of $\xi$ and common to
all $k$.  Suppose its values on these cells are bounded measurable
functions of $\xi$, and that
\begin{equation}\label{eq:frequency-continuous-condition}
 |\mu_{a,b,k}(\xi)|\leq C_1,
 \qquad
 \sum_k\int_{\Xi_D}|\mu_{a,b,k}(\xi)|^2\,\dd\gamma(\xi)
 \leq C_1^2.
\end{equation}
Define
\begin{equation}\label{eq:frequency-continuous-kernel}
 \mathcal K_{a,b}(t)
   =\int_{\Xi_D}\sum_{k\in\mathcal I}
          \mu_{a,b,k}(\xi)\Dil_{L_k}H_\xi(t)\,\dd\gamma(\xi).
\end{equation}
Then, for complex $G_v\in L^3(\R^2)$,
\begin{equation}\label{eq:frequency-continuous-estimate}
 |\Lambda_{\mathcal K}(G_0,G_1,G_2)|
       \leq C(q,C_1)\prod_{v=0}^2\|G_v\|_3.
\end{equation}
\end{corollary}

\begin{proof}
First take smooth compactly supported $G_v$.  Within any fixed
step choice, the kernel in
\eqref{eq:frequency-continuous-kernel} is continuous in $t$:
the frequency block is bounded, the coefficients are bounded and
measurable, and the Gaussian formulas supply a common majorant
for dominated convergence.  After performing the frequency
integration, the trilinear integral is therefore the limit of
its lattice Riemann sums.  The finitely many rectangular step
boundaries have measure zero.  The mesh condition
\eqref{eq:frequency-mesh} holds for all sufficiently fine lattices,
and the discrete norms converge to $\|G_v\|_3$.
Proposition~\ref{prop:frequency-block} proves
\eqref{eq:frequency-continuous-estimate} for these inputs.

To pass to general $L^3$ inputs, note that for the fixed finite
block and scale set there is an integrable majorant $k_*(t)$
with $|\mathcal K_{a,b}(t)|\leq k_*(t)$, uniformly in $(a,b)$.
One may take a constant times the finite sum of the dilates of
$\sup_{\xi\in\Xi_D}|H_\xi|$, which has Gaussian decay.
For every integrable nonnegative $k$,
\[
 \iiint |G_0(a,b)G_1(b,c)G_2(c,a)|k(a+b+c)
                         \,\dd a\,\dd b\,\dd c
 \leq\|k\|_1\prod_v\|G_v\|_3.
\]
Indeed, set $t=a+b+c$, apply H\"older's inequality in $(a,b)$
for fixed $t$, and use the determinant-one changes of variables
for the other two edge functions.  Thus the defining integral
is absolutely convergent and continuous in the three $L^3$
inputs.  Smooth approximation completes the proof.  The final
constant comes from the Proposition and remains independent of
$D$, the number of scales, and the particular step choices.
\end{proof}
